\chapter{The Primal Assumptions}
\label{chap:primal-assumptions}

\section{The Quest for the Presumptionless Theory}

The holy grail of theoretical physics and cosmology has always been to find a singular,
self-evident first principle from which all of reality can be derived.

Our intuition seems to \textbf{demand} a framework that does not smuggle in arbitrary starting conditions.

We are dissatisfied with the Standard Model of particle physics precisely because it forces us to accept, without explanation, at least 19 free parameters: fine-tuned constants such as the fine-structure constant $\alpha$, the gravitational constant $G$, and the masses of various particles.
To a philosophically minded software developer, accepting these numbers as "just because" feels like intellectual surrender.

And software developers are apparently not the only ones dissatisfied with mainstream physics.
Several theories seek to derive the laws of physics from minimal assumptions.

Ludwig Boltzmann proposed that in a system at thermal equilibrium (maximum entropy), rare fluctuations can temporarily produce ordered structures.
In such a framework, sufficiently complex configurations—including observers—can arise given enough time. 

Max Tegmark's Mathematical Universe Hypothesis is another rival.
It posits that all mathematically consistent structures physically exist, requiring no external creator or physical spark.

A third major candidate is Archibald Wheeler's "It from Bit" paradigm, where physical reality is derived entirely from the binary processing of quantum information.

However, none of these theories are truly presumptionless. Bolzmann assumes time and dynamics.
Tegmark assumes that mathematical structures exist as a fundamental substrate, and 'It from Bit' takes the validity of quantum mechanics for granted.
In all cases, the 'minimal' starting point is actually quite a lot to assume.


\section{The Barren Landscape of Lawless Theories}

The fundamental trouble with these lawless, ultimate first-principle theories is that they do not yield the universe we actually observe.
They struggle to get even two planets to orbit each other with inertia, let alone derive the smooth, continuous spacetime metric.

It would actually make sense to assume that reality is born from pure probabilistic fluctuations in an infinite sea of chaos.
It would then be something we could truly understand. However, we immediately run into the Boltzmann Brain paradox.
In a truly random infinity, it is statistically far more probable for a single conscious brain with false memories to fluctuate into existence than a vast, ordered universe governed by rigid laws. 

If the universe is based on mathematics, we are immediately plagued by the measure problem.
There are infinitely more chaotic, non-computable, and glitchy mathematical structures than there are simple, elegant ones.
Without adding an external, ad-hoc rule favoring simplicity, this theory predicts we should live in a wildly unstable, incomprehensible reality. 

Algorithmic Information Theory is an exceptional tool for ledger-keeping and measuring quantum states, but it fails to natively explain why mass resists acceleration (inertia) or how smooth, continuous gravity emerges from discrete bits.


\section{The Inevitability of the Starting Assumption}

If one seeks to minimize unexplained assumptions as far as possible, one eventually arrives at the statement:
\textit{This theory assumes nothing}, which is already an assumption.

This reflects, in a broader philosophical sense, a consequence reminiscent of G\"odel's incompleteness:
every sufficiently rich formal system appears to require at least some starting point that cannot be derived entirely from within the system itself.
At bare minimum, there must exist at least \textit{one} foundational premise.

And if a system requires at least one assumption in order to exist, what dictates that there must be only one? Why not two? Why not nineteen?

The nineteen seemingly arbitrary constants of the Standard Model, the specific geometry of quantum Hilbert spaces, and the fundamental structure of General Relativity may represent precisely such irreducible starting assumptions: a minimal set of brute facts required for this universe to exist at all.

The remarkable success of these theories may point toward a deeper and more uncomfortable truth about the nature of reality and logic itself.
Reductionism may have a floor. We can reduce chemistry to physics, and physics to quantum fields,
yet we may ultimately arrive at a bedrock of predefined constants and structural axioms that cannot be reduced any further.

The goal of a theory of everything may therefore not be to eliminate all assumptions,
but to determine the smallest and most coherent set of assumptions from which everything else necessarily follows.

